\documentclass{article}
\usepackage{amsmath,amssymb}
\usepackage{xurl}
\usepackage[hidelinks]{hyperref}
% Shared commands for notes in docs/paper-gaps/.

\DeclareUnicodeCharacter{2080}{\ifmmode{}_{0}\else\textsubscript{0}\fi}
\DeclareUnicodeCharacter{2081}{\ifmmode{}_{1}\else\textsubscript{1}\fi}
\DeclareUnicodeCharacter{2082}{\ifmmode{}_{2}\else\textsubscript{2}\fi}
\DeclareUnicodeCharacter{2099}{\ifmmode{}_{n}\else\textsubscript{n}\fi}

\newcommand{\Lean}{\textsc{Lean}}
\ExplSyntaxOn
\NewDocumentCommand{\leanid}{m}{
  \ifmmode
    \textsf{\detokenize{#1}}
  \else
    \group_begin:
    \urlstyle{sf}
    \tl_set:Nn \l_tmpa_tl {#1}
    \tl_replace_all:Nnn \l_tmpa_tl {\_} {_}
    \exp_args:NV \path \l_tmpa_tl
    \group_end:
  \fi
}
\ExplSyntaxOff
\providecommand{\tr}{\operatorname{tr}}
\newcommand{\tnleanIssueBase}{https://github.com/LionSR/TNLean/issues/}
\newcommand{\tnleanPrBase}{https://github.com/LionSR/TNLean/pull/}
\newcommand{\tnleanDocBase}{https://sirui-lu.com/TNLean/docs/}
\newcommand{\ghissue}[1]{\href{\tnleanIssueBase#1}{GitHub issue \##1}}
\newcommand{\ghpr}[1]{\href{\tnleanPrBase#1}{PR \##1}}
\newcommand{\doclink}[1]{\href{\tnleanDocBase#1.html}{\path{#1}}}

% Dirac notation and the identity operator, as in blueprint/src/macros/common.tex.
% \providecommand keeps note-local definitions authoritative.
\usepackage{dsfont}
\providecommand{\ket}[1]{|#1\rangle}
\providecommand{\bra}[1]{\langle #1|}
\providecommand{\braket}[2]{\langle #1 | #2 \rangle}
\providecommand{\ketbra}[2]{|#1\rangle\!\langle #2|}
\providecommand{\Id}{\mathds{1}}
\providecommand{\one}{\mathds{1}}
\providecommand{\C}{\mathbb{C}}
\providecommand{\MN}[1]{M_{#1}(\C)}
\providecommand{\MD}{M_D(\C)}

% Verdict marker: \gapnote{<kind>}{<status>}. Kind is the policy
% classification (clarification, local-correction, scope-restriction,
% unfaithful, false-source, open-gap); status is open, wip (elimination
% underway), resolved, or historical. Machine-read by the site index;
% prints a verdict line.
\newcommand{\gapnote}[2]{\par\noindent\textbf{Verdict:} #1 (#2).\par}

\title{The Three-Object Ising Twist: Channels and Formal Scope}
\author{}
\date{2026-10-02}
\begin{document}
\maketitle
\gapnote{scope-restriction}{open}

\section{The construction and its four channels}

The Ising fusion rules and $F$-symbols are given in
Ref.~\cite{Bultinck2017Anyons}, Appendix D.2. The weighted bond-object
construction is a project example motivated by the question following
Theorem 4.14 of Ref.~\cite{Cirac2016MPDO_arXiv}; it is not a tensor printed
in either paper. Its local construction record is
\path{Notes/OpenProblemsTN/checks/asym_ising_action_data.md}, sections
0 and 1.4. The relevant equations are stated here independently of that
local record.

Let $A_1,A_\psi,A_\sigma$ be the three Ising operator tensors, of auxiliary
dimensions $3,3,4$. Write $\widehat A_\sigma=\sqrt2 A_\sigma$ and
$\Theta_3=5A_1\oplus3A_\psi\oplus2A_\sigma$. On a positive-length ring,
$O_L(A_1)O_L(\widehat A_\sigma)=O_L(\widehat A_\sigma)$,
$O_L(A_\psi)O_L(\widehat A_\sigma)=O_L(\widehat A_\sigma)$, and
$O_L(\widehat A_\sigma)^2=2^L(O_L(A_1)+O_L(A_\psi))$. Hence
\begin{align}
 O_L(\Theta_3)O_L(\widehat A_\sigma)
 &= (5^L+3^L)O_L(\widehat A_\sigma)
  +(2\sqrt2)^L\bigl(O_L(A_1)+O_L(A_\psi)\bigr).
 \label{eq:ising_three_fusion}
\end{align}
The same formula holds for traces along every nonempty word. It supplies
four weighted normal blocks, namely $5\widehat A_\sigma$,
$3\widehat A_\sigma$, $2\sqrt2 A_1$, and $2\sqrt2 A_\psi$.
Asymmetric compression gives $40=4+4+3+3+26$, with twenty-six
one-dimensional zero diagonal blocks in an upper-triangular presentation.
An explicit change of auxiliary coordinates strengthens this to a block-diagonal
presentation. It is the direct sum of the twenty-four-dimensional gauge for
the two abelian bond objects and the sixteen-dimensional gauge for
$\sigma\otimes\sigma$. The inverse of the latter is half its transpose.
After separating the four active blocks from the twenty-six zero coordinates,
the off-diagonal remainder vanishes. The left projection and right inclusion
intertwine every letter with its weighted target. This statement concerns
this explicit choice; the compression obtained only from the trace theorem
need not be the same choice.

The three-term coefficient $5^L+3^L+2^L$ for a single $\sigma$ target,
and the three-slot compression requested in issue 7833, are not the
construction's fusion rule. The last bond object contributes to $1$ and
$\psi$ through $\sigma\times\sigma=1+\psi$. For example, at the one-letter
word $(h,h')=(0,0)$, the trace of the stacked tensor is $2\sqrt2$, whereas
the trace of $\widehat A_\sigma^{00}$ is zero. This already excludes the
proposed single-target formula.

\section{The full boundary tensors remain distinct}

At the numerical point $(\lambda_1,\lambda_\psi,\lambda_\sigma)
=(1/2,3/10,1/5)$, the unrescaled bond object is $\Theta_3/10$.
This scalar relation concerns the local forty-dimensional product above.
The full round-45B boundary tensor of flag $f$ has site space
$\mathbb C^4\otimes P_f\otimes\mathbb C^4$, of dimension
$48,48,64$ for $f=1,\psi,\sigma$, and horizontal indices
$(h,r,s)$ with $h\in\{0,\ldots,9\}$ and $r,s\in\{0,\ldots,3\}$.
Its matrices take the form
\begin{align}
 M_f^{(h,r,s),(h',r',s')}
 &=X_f^{h,r,s,h'}\otimes E_{r's'}.
 \label{eq:ising_full_boundary}
\end{align}
Their fusion is the separate assertion
\begin{align}
 O_L(M_f)O_L(M_{f'})
 &=\sum_{b,c}\lambda_b^L N^c_{fbf'}O_L(M_c),
 \label{eq:ising_full_boundary_fusion}
\end{align}
where $N^c_{fbf'}$ is the multiplicity of $c$ in $f\otimes b\otimes f'$.
In particular $N^\sigma_{\sigma\sigma\sigma}=2$.
Neither (\ref{eq:ising_three_fusion}) nor a common scalar rescaling proves
(\ref{eq:ising_full_boundary_fusion}).

At $(f,f')=(\sigma,\sigma)$ the stacked auxiliary dimension is $4096$,
the alphabet has $25600$ letters, and the active targets have dimensions
$48,48,48,48,64,64$. Direct exhaustive kernel evaluation of dense matrices
at this size is not a suitable extension of the existing small examples.
No exact full-tensor certificate is supplied here. The elimination plan is
to define the reduced letters $X_f$ at the rational point, prove their
fusion from the sparse Ising $F$-move identities and the register matrix-unit
contraction, and then restore the outer register in
(\ref{eq:ising_full_boundary}). Positivity, normalization and refinement
channels must be established for those same literal tensors.
This remaining scope is the second part of
\url{https://github.com/LionSR/TNLean/issues/7833}; it must not be marked
formalized on the basis of the forty-dimensional compression.

\section{Formal declarations}

The local result is in
\path{TNLean/MPS/Examples/Ising/IsingThreeObjectTwist.lean}:
\leanid{IsingTwist.thetaThree_mpo_mul_isingSigma},
\leanid{IsingTwist.isingBondObjectThree_trace_evalWord},
\leanid{IsingTwist.isingBondObjectThree_compression_exists}, and
\leanid{IsingTwist.isingBondObjectThree_z_eq}.
The explicit split compression is in
\path{TNLean/MPS/Examples/Ising/IsingThreeObjectGauge.lean}:
\leanid{IsingTwist.isingThreeSplitCompression},
\leanid{IsingTwist.isingBondObjectThree_split_remainder},
\leanid{IsingTwist.isingBondObjectThree_mul_right_eq_right_mul}, and
\leanid{IsingTwist.isingBondObjectThree_left_mul_eq_mul_left}.

\bibliographystyle{alpha}
\bibliography{references}
\end{document}
